\documentclass[11pt,a4paper]{article}
\usepackage[T1]{fontenc}
\usepackage[utf8]{inputenc}
\usepackage{lmodern}
\usepackage[margin=25mm]{geometry}
\usepackage{amsmath,amssymb,mathtools}
\usepackage{booktabs,array}
\usepackage[hidelinks]{hyperref}
\usepackage{microtype}
\setlength{\parskip}{0.2em}
\setlength{\parindent}{1.2em}
\newcommand{\diag}{\operatorname{diag}}
\newcommand{\norm}[1]{\left\lVert#1\right\rVert}
\newcommand{\code}[1]{\texttt{#1}}
\title{Experimental Design, Learning-Rate Selection and Batch Execution\\[0.35em]
\large AISC5020: Module 3 Report Sections (Sec.~4 and Sec.~5.1--5.2)}
\author{CHEN Yuchan}
\date{19 October 2026 (reviewed revision)}
\begin{document}
\maketitle
\begin{abstract}
These sections specify how the paired experiments are generated, tuned and
executed. Every problem, rotation and initial point is rebuilt from its own seed
key, so pilot, tuning and evaluation instances never overlap and any single run
can be reproduced in isolation. Learning rates are chosen from sixteen
log-spaced candidates under two regimes, transfer and per-orientation
retuning, with gradient-descent and heavy-ball candidates kept strictly inside
their stability intervals. A pilot fixes the candidate ranges before any
evaluation run. Fixed-budget accuracy, threshold attainment and three failure
categories are recorded separately. The historical Release dataset is
distinguished from subsequent implementation validation. These sections do
not rank optimizer performance.
\end{abstract}

\noindent\emph{Integration note.} Section numbers follow the proposal, so
Sections~\ref{sec:design} and~\ref{sec:tuning} can be pasted into the group
report. Section~5.3--5.4 (outcome measures and uncertainty) belongs to Module~4;
the quantities they use are defined here only as far as the run records require.
Section~\ref{sec:implementation} is supporting material for the appendix.

\setcounter{section}{3}
\section{Experimental design}\label{sec:design}

\subsection{Problem families and paired instances}
All problems are the quadratics $f_Q(w)=\tfrac12 w^\top QDQ^\top w$ of
Section~1 with $L=1$ and $\mu=1/\kappa$, where the eigenvalues are
\begin{equation}
  d_i=\kappa^{-(p-i)/(p-1)},\qquad i=1,\dots,p,
\end{equation}
so $D=\diag(1/\kappa,1)$ when $p=2$ and the spectrum is log-spaced otherwise.
We use $\kappa\in\{10,10^2,10^3,10^4\}$ and, as an isotropic control,
$\kappa=1$ (Section~\ref{sec:controls}).

For each Gaussian direction $u$ we set $z_0=\sqrt{2/(u^\top Du)}\,u$ and
$w_0=Qz_0$, so $f_I(z_0)=f_Q(w_0)=1$ and $\norm{w_0}=\norm{z_0}$ (eq.~(3)).
A run on $(Q,w_0)$ is always compared with the run on $(I,z_0)$ that starts from
the same canonical point; the latter is called its \emph{diagonal reference}.
Every algorithm uses the same instances.

\paragraph{Two-dimensional study.} For each $\kappa$ we draw ten initial
directions and share them across the angles
$\theta\in\{0^\circ,15^\circ,\dots,90^\circ\}$, with $Q=Q_\theta$ from eq.~(25).
The $0^\circ$ run is the reference for every angle and is executed once; it is
not counted again for each angle. The $90^\circ$ rotation is a signed
permutation and serves as a control for which Adam correspondence is exact.

\paragraph{Higher-dimensional study.} For $p=20$ and each $\kappa$ we draw ten
Haar rotations and, for each rotation, five initial directions. Each of the 50
rotated runs has its own diagonal reference with the same $z_0$.

\paragraph{Algorithms.} The primary comparison contains gradient descent,
heavy-ball momentum with $\beta=0.9$, Adam with $(\beta_1,\beta_2,\epsilon)=
(0.9,0.999,10^{-8})$ and common-eigenbasis Adam with the generator basis; the
diagonally preconditioned gradient method is a secondary comparison. All
methods are called through one interface, so each gradient evaluation counts
as one call and loss bookkeeping is not charged to the budget.

\subsection{Randomization and reproducibility}
Each random object has a dedicated stream
\code{SeedSequence(root, key)}, where the key records the stage (pilot, tuning
or evaluation), the purpose (rotation or initial direction), the study,
$\kappa$, the tuning regime, the pilot round, the orientation and the index.
Streams from different stages are therefore disjoint by construction, a fresh
pilot round uses fresh seeds, and any run can be rebuilt from the key stored in
its record without replaying the rest of the study. Haar rotations come from
Gaussian QR factorizations with diagonal-sign correction. All arithmetic is in
double precision. Each record stores the key of its rotation and initial
direction, the configuration and a hash of its specification. The reviewed
execution format additionally records source provenance before running; it
does not infer old execution provenance from the files present at report time.

\subsection{Controls}\label{sec:controls}
With $\kappa=1$ the Hessian is the identity for every $Q$, so any difference
between a rotated run and its reference reflects only the coordinate
representation of the initial point. The $0^\circ$ and $90^\circ$ angles are
signed permutations, for which Adam trajectories correspond exactly.
Gradient descent and heavy ball are rotation-equivariant (Results P1 and P2),
so under a shared learning rate their paired differences provide an end-to-end
check of the pairing itself. Selected runs with unit-norm initial points
($\norm{z_0}=1$) or scaled copies can reuse the same underlying directions,
but their normalization and scale are part of the study definition. Such
controls must be specified before freezing a new output directory.

\subsection{Run records and failure categories}\label{sec:records}
Every evaluation run has a budget of $50{,}000$ gradient calls and records
\begin{equation}
  r_B=\frac{f(w_B)}{f(w_0)},\qquad B\in\{10^3,10^4,5\times 10^4\},
\end{equation}
together with the first attainment time $T=\min\{t:r_t\le 10^{-6}\}$, checked at
every step, whether $r_t\le10^{-6}$ still holds at the final step, and the
largest recorded $r_t$ after $T$. Runs are never stopped at the threshold. A
run ends in one of four states: \code{finite\_cap} (budget completed),
\code{explosive} ($r_t>10^{12}$), \code{nonfinite} (an infinite or undefined
value in the iterate, gradient or optimizer state), or
\code{not\_run\_tuning\_failure} when every candidate rate failed during
tuning. Explosive and non-finite runs keep their last finite iterate, the
failing step and their trace. No run is assigned an artificial attainment time
and no failed run is removed from a denominator. For each rotated run the
paired record lists both outcomes, the four-way attainment category, the
attainment ratio when both runs succeed, and $\Delta_B$ of eq.~(32) at each
checkpoint where both values are finite. It also states whether the pair
coincides in exact arithmetic (gradient descent, heavy ball and
common-eigenbasis Adam under shared settings; Adam and diagonal
preconditioning only for signed permutations) and the first recorded step at
which the two traces of $\log_{10} r_t$ differ by more than $10^{-6}$. These
pairs provide qualitative diagnostics of numerical sensitivity. Their
distribution is not automatically a calibrated error distribution for methods
with different learning rates or different adaptive states.

\section{Tuning and evaluation}\label{sec:tuning}

\subsection{Learning-rate selection}
Every algorithm has sixteen log-spaced candidate rates. A candidate has a
budget of $5{,}000$ gradient calls from each tuning initialization, without
threshold-based early stopping, and receives the score
\begin{equation}
  S(\eta)=\operatorname{median}\Big[\log_{10}\max\{r_{5000},10^{-16}\}\Big].
\end{equation}
A candidate with any explosive or non-finite tuning run receives the worst
score; a run that only misses $10^{-6}$ within the budget is not a failure.
The candidate with the smallest score is selected, ties go to the smaller rate,
and if every candidate fails the cell is recorded as a tuning failure rather
than given a rate.

Candidate ranges for gradient descent and heavy ball lie strictly inside the
stability intervals of Section~1: the largest candidate is $0.999$ times the
upper limit, that is $\eta\le1.998$ for gradient descent and $\eta\le3.7962$ for
heavy ball with $\beta=0.9$, whose limit is $2(1+\beta)/L=3.8$. The Adam
variants and the preconditioned method use logarithmic ranges without an
analytical cap.

\paragraph{Transfer regime.} For each algorithm, dimension and spectrum, one
rate is selected from four tuning orientations, the diagonal basis and three
random orientations, with five tuning initializations each. In two dimensions
the random orientations are angles drawn uniformly from $(0^\circ,90^\circ)$,
so they generally differ from the evaluation grid. The rate is then fixed for
every evaluation orientation.

\paragraph{Retuning regime.} For each evaluation orientation, including the
diagonal reference, a rate is selected from five tuning initializations of that
same orientation. These initializations come from the tuning stream, so the
rate is evaluated on fresh initial points.

\paragraph{Spectral references.} Gradient descent with $\eta^\ast=2/(L+\mu)$
and heavy ball with $(\eta^\ast_{\mathrm{HB}},\beta^\ast_{\mathrm{HB}})$ from
eqs.~(7) and~(10) are run on the same evaluation instances. They use knowledge
of $\mu$ and $L$ and are reported separately from the equal-budget comparison.

\paragraph{Horizons.} Rates are selected at $5{,}000$ calls and evaluated at
$50{,}000$. A rate that is best at the shorter horizon need not be best at the
longer one; the gradient calls spent on the pilot, on each tuning regime and on
evaluation are reported separately.

\subsection{Pilot and frozen protocol}
The pilot uses seeds that are never used again. For every algorithm, study and
spectrum it runs the full candidate grid on two orientations (including the
diagonal basis) with three initializations each, and applies the selection
rule above. If any selected rate is a boundary optimum, the range of that
algorithm is widened by a factor of ten on that side and the pilot is repeated
with fresh seeds, at most twice. An upper boundary never moves past
$0.999$ times a stability limit; an optimum there is reported as capped by
stability rather than expanded. A boundary that remains after the last round
is reported as a range limitation. The reviewed format freezes the candidate
rates and the study definitions, including dimension, spectra, normalization,
scale, orientation schedule and sample counts. Evaluation budgets and the other
configuration fields also remain fixed within that output directory. The only
permitted reduction is to omit entire studies whose remaining definitions are
unchanged; modifying a study requires a new pilot in a new directory. The
original historical format did not enforce all of these distribution checks.

The $10^{-16}$ floor can make neighbouring scores equal. Under the default
smaller-rate tie rule, the pilot counts a selected edge candidate as a boundary
optimum only when its score is strictly better than that of its neighbour; a tied edge is recorded as a plateau
and is not expanded. This is an operational pilot rule, not a proof that a
finite grid locates the global optimum. If the scores actually attain the
$-16$ floor, no candidate can improve this floored criterion; that statement
does not extend to unrounded losses or the longer evaluation horizon.

\paragraph{Pilot outcome.} The pilot used 5{,}760 runs ($2.85\times10^{7}$
gradient calls) on pilot seeds only. Table~\ref{tab:pilot} lists the frozen
ranges. Gradient descent and heavy ball kept their initial ranges: their optima
at $\kappa=10^4$ (and at $\kappa=10^3$ for two-dimensional gradient descent)
lie at the stability-capped upper end. For heavy ball at $\kappa\in\{1,10\}$
all sixteen candidates reach the floor, which is the plateau case above. Adam
had a strict upper-boundary optimum in the two-dimensional $\kappa=10^4$ cell,
so its upper end was raised to $10$; in the second round, with fresh seeds,
every optimum was interior. Adam and common-eigenbasis Adam therefore use grids
of different spacing, which must be kept in mind when the two are compared.

\begin{table}[h]
\centering
\caption{Candidate ranges after the pilot (sixteen log-spaced rates each).}
\label{tab:pilot}
\small
\begin{tabular}{@{}lllll@{}}
\toprule
Method & Initial range & Rounds & Frozen range & Boundary note \\
\midrule
Gradient descent & $[0.002,\,1.998]$ & 1 & unchanged & capped by stability \\
Heavy ball, $\beta=0.9$ & $[0.0038,\,3.7962]$ & 1 & unchanged & capped by stability; plateau \\
Adam & $[10^{-4},\,1]$ & 2 & $[10^{-4},\,10]$ & upper end expanded once \\
Eigenbasis Adam & $[10^{-4},\,1]$ & 1 & unchanged & none \\
Diagonal preconditioning & $[0.002,\,2]$ & 1 & unchanged & none \\
\bottomrule
\end{tabular}
\end{table}

Under the tie rule, a floor plateau selects the smallest candidate among
the tied floored scores. Those candidates have the same recorded tuning score,
but need not have the same raw loss, first-attainment time or later accuracy.
The rule does not minimize first-attainment time, so comparisons at small
$\kappa$ should be read with this in mind. An
alternative that breaks score ties by the median attainment time is available
as a configuration option; changing it requires a new pilot.

\section{Implementation and verification}\label{sec:implementation}
The protocol is implemented in four files. \code{part3\_design} builds the
spectra, rotations, initial points and run units; \code{part3\_tuning} holds the
candidate grids, stability limits, scoring and pilot rules;
\code{part3\_experiments} runs the pilot, tuning, evaluation and reporting
stages in parallel worker processes and resumes interrupted stages from saved
records; and \code{verify\_part3} checks the saved outputs with separate code.
The original verifier rebuilt sampled instances from stored seed keys without
using the design module and re-executed them through the shared runner. It
recomputed the selected rates from raw tuning records and checked seed-stage
isolation, canonical-point pairing and shared transfer rates. The reviewed
verifier adds completeness and provenance checks; its current verification
scope is reported separately from the historical main-study numbers below.
The complete expected run plan, candidate records, selections, pair and CSV
fields, and budget totals are checked independently of the numerical sample.
Checkpoint comparisons use explicit tolerances (default relative tolerance
$10^{-8}$ and absolute tolerance $10^{-12}$), not a cross-platform bitwise
identity promise. Passing requires both the structural and sampled numerical
checks, reported separately as \code{structural\_checks\_passed} and
\code{numerical\_checks\_passed}. A failed writable verification replaces an old passed marker with a
failed result; \code{--read-only} writes no verification file or manifest.

\paragraph{Execution provenance and safe resumption.} Before a new pilot,
format version 2 creates \path{execution_origin.json} by exclusive creation.
It records the initial configuration, source-file hashes normalized from CRLF
to LF, a source fingerprint and the execution environment. Each new raw record
carries \code{execution\_source\_sha256}. Subsequent stages check the origin
instead of replacing it with the current source identity; source changes or
changed Python/NumPy versions require a fresh directory. The manifest's
execution identity derives from this origin. A rejected guard does not update
the stored run configuration or manifest, and preserves the previous
verification marker. After preflight succeeds, a new execution stage clears
that marker so a later stage cannot inherit an old passed claim.

An old output directory without this origin is historical evidence, not a
cache to silently upgrade. The execution stages refuse to append to or
relabel it; a read-only verifier may inspect it without inventing version-2
provenance or requiring its historical source to equal the current checkout.
New report generation validates completeness against the active
run plan and filters compatible whole-study subsets, rather than exporting
whatever records happen to exist. The revised budget summary records
\code{active\_studies} and \code{report\_scope}, making a subset's counts
explicit rather than mixing in omitted-study caches.
Resumption removes only an unfinished final
JSONL fragment without a line ending before appending; malformed complete
lines or interior records are rejected. A valid last record without a newline
is retained and separated from the next record.
Seeds, dimensions and count/budget fields require integers, not Booleans or
silently truncated fractions. Condition numbers retain the original seed-key
restriction to finite integer-valued numbers at least one (both $10$ and
$10.0$ are accepted); this is an implementation namespace restriction, not a
mathematical restriction on quadratics. Regression checks for these cases complement
the original generator, selection, spectral-reference and end-to-end tests.

\paragraph{Validation of the reviewed implementation.} The final integrated
check passed all 89 tests, including 44 Module-3 tests (18 workflow,
11 provenance/recovery and 15 verifier tests), and all default experiment
checks. The Module-3 smoke check passed the full structure of its 168 pilot,
600 tuning and 240 evaluation records, with 60 numerical replays per stage.
These checks do not rerun the historical main study.

\paragraph{Historical-data replay and its limit.} A separate read-only
check passed the full structure of the Release dataset: $5{,}760$ pilot,
$52{,}000$ tuning and $10{,}200$ evaluation records, with $6{,}600$ pairs.
Its numerical replay sampled 150 runs per stage, 450 in total, and did
\emph{not} pass at relative tolerance $10^{-8}$ and absolute tolerance
$10^{-12}$. There were 215 failed field comparisons across 56 runs, not
215 failed runs: 7 pilot, 11 tuning and 38 evaluation runs. They comprised
23 Adam and 33 eigenbasis-Adam runs; the other sampled runs passed. One pilot
first-attainment comparison also failed. The Adam,
eigenbasis-Adam, runner and problem kernels were unchanged from \code{575ebf4}.
The 450-run check preceded the final CSV-metadata-only verifier tightening.
The final checker separately passed the same full structural checks and one
numerical replay per stage in read-only mode. This three-run numerical check
does not supersede the failed 450-run check. These are separate checks, not
one all-passing final-checker command; the
\href{https://github.com/kvvvvvv23/aisc5020-coordinate-geometry/blob/main/docs/part3_validation_2026-10-19.json}{audit summary (JSON)}
records their scopes and source identities.

In a controlled single-run rerun, changing only the first component of $w_0$
by one unit in the last place (perturbation norm $2.78\times10^{-17}$) shifted
the final relative loss from $9.46\times10^{-5}$ to $1.08\times10^{-7}$ and
changed final threshold retention; this example's first attainment remained
at step 122. This supports numerical sensitivity, but
does not identify the specific cause of the historical discrepancies. The
released data are therefore not fully reproduced numerically in the current
environment, even though their structure is verified.

\paragraph{Historical Release execution.} The original
\href{https://github.com/akaComic/aisc5020-coordinate-geometry/releases/tag/part3-v1}{\code{part3-v1} Release}
records the run made on 18--19 October 2026 at source commit
\code{575ebf4}, using Python 3.13.16 and NumPy 2.5.4 on Linux, two workers
and single-thread BLAS settings. It produced
$52{,}000$ tuning runs and $10{,}200$ evaluation runs. Every tuning cell
selected a rate; the 47 boundary optima were all gradient-descent or heavy-ball
optima at the stability cap. All records that did not end in
\code{finite\_cap} were large-rate candidates of the diagonal method during
pilot and tuning; every evaluation run completed its budget and every pair
attained the threshold. The Release contains a passed verification record with
150 sampled runs per stage; its reported checks include all 500 selections.
These are historical results, not a claim that the full main experiment was
rerun after the review fixes.

A separate read-only review checked the published archive's SHA256 against
the GitHub asset digest and matched all 19 source-file hashes to the historical
snapshot (with line endings represented as LF). It independently checked
the initial-seed and metadata pairing of all $6{,}600$ pairs and recomputed
their checkpoint $\Delta_B$ values; the largest numerical difference was
$1.78\times10^{-15}$. This corroborates the released tables and provenance.
That archive-integrity check did not rerun trajectories; the separate
450-run numerical replay and its failures are described above.

\paragraph{Observed late-time separation.} First consider only pairs whose
required settings are shared (the field \path{correspondence_expected} is true).
In the historical dataset there are $4{,}460$ such pairs. None separates at a
recorded point before both runs attain $10^{-6}$; the gradient-descent,
heavy-ball and diagonal-method pairs in this subset do not separate at all.
This statement does not cover all retuned pairs, since different selected
learning rates can themselves change the trajectories. Among the expected
correspondence pairs, $1{,}118$ Adam or common-eigenbasis-Adam pairs separate
later, with median recorded separation step $17{,}000$. Their late
$\Delta_B$ can be large even though the idealized trajectories coincide.

Of the $1{,}700$ ordinary-Adam evaluation runs, 678 were observed above
$10^{-6}$ at a saved point after first attainment, and 128 ended above it.
The rebound count is based on records spaced $1{,}000$ updates apart; it is
not a complete count of all between-record threshold crossings. In the
stored expected-correspondence pairs the first-attainment times are identical and
the absolute $10^3$-call paired log gap is at most $3\times10^{-10}$.
These observations describe the saved dataset, not equality across environments.

\paragraph{A conditional local model for the spikes.} A decaying second
moment can increase the effective update scale, but Adam's first momentum
must also be retained. To see the distinction from GD, consider a
one-dimensional quadratic with curvature $\lambda>0$, neglect late-time bias
corrections, and temporarily freeze
$c=\sqrt{\bar v}+\epsilon>0$. Write $\gamma=\eta/c$. The resulting linear model is
\[
 m_{t+1}=\beta_1m_t+(1-\beta_1)\lambda x_t,\qquad
 x_{t+1}=x_t-\gamma m_{t+1}.
\]
Eliminating the first moment gives
\begin{equation}
 x_{t+1}=\bigl[1+\beta_1-\gamma(1-\beta_1)\lambda\bigr]x_t
          -\beta_1x_{t-1}.
\end{equation}
For $0\leq\beta_1<1$, the strict Schur condition for this frozen model is
\begin{equation}
 0<\gamma\lambda<\frac{2(1+\beta_1)}{1-\beta_1}.
 \label{eq:frozen-adam}
\end{equation}
Thus the upper bound is $38$ when $\beta_1=0.9$, rather than the GD value $2$.
For example, at $\gamma\lambda=3$ the roots are approximately
$0.8\pm0.509902\,\mathrm{i}$, with modulus $\sqrt{0.9}<1$;
at $38$ the roots are $-1$ and $-0.9$, so the boundary is not strictly stable.
The actual denominator is $\sqrt{\hat v_t}+\epsilon$, so its effective scale
is bounded above by $\eta/\epsilon$ and cannot be assumed to grow without limit.

Equation~\eqref{eq:frozen-adam} is a one-dimensional, frozen-denominator,
late-time approximation, not a global stability theorem for adaptive Adam.
When gradient squares contribute little to the second-moment update, its
decay may help explain a subsequent rise in effective scale. The complete
adaptive dynamics, changing moments and finite precision still have to be
considered. A double-precision rerun alone does not prove that the same
observed spikes occur in exact arithmetic, nor does it establish mathematical
chaos. The recorded loss traces show late rebound and sensitivity; they do not
supply such a proof.

Late-checkpoint gaps for the Adam variants can therefore reflect numerical
sensitivity as well as coordinate or tuning effects. Matched eigenbasis-Adam
and signed-permutation controls help diagnose that sensitivity, but are not
an error distribution to subtract from every other method. The observed
floating-point separation does not contradict Result P4's exact-arithmetic
common-basis identity.

\end{document}
